% Options for packages loaded elsewhere
\PassOptionsToPackage{unicode}{hyperref}
\PassOptionsToPackage{hyphens}{url}
\documentclass[
]{article}
\usepackage{xcolor}
\usepackage{amsmath,amssymb}
\setcounter{secnumdepth}{5}
\usepackage{iftex}
\ifPDFTeX
  \usepackage[T1]{fontenc}
  \usepackage[utf8]{inputenc}
  \usepackage{textcomp} % provide euro and other symbols
\else % if luatex or xetex
  \usepackage{unicode-math} % this also loads fontspec
  \defaultfontfeatures{Scale=MatchLowercase}
  \defaultfontfeatures[\rmfamily]{Ligatures=TeX,Scale=1}
\fi
% \usepackage{lmodern}
\ifPDFTeX\else
  % xetex/luatex font selection
\fi
% Use upquote if available, for straight quotes in verbatim environments
\IfFileExists{upquote.sty}{\usepackage{upquote}}{}
\IfFileExists{microtype.sty}{% use microtype if available
  \usepackage[]{microtype}
  \UseMicrotypeSet[protrusion]{basicmath} % disable protrusion for tt fonts
}{}
\makeatletter
\@ifundefined{KOMAClassName}{% if non-KOMA class
  \IfFileExists{parskip.sty}{%
    \usepackage{parskip}
  }{% else
    \setlength{\parindent}{0pt}
    \setlength{\parskip}{6pt plus 2pt minus 1pt}}
}{% if KOMA class
  \KOMAoptions{parskip=half}}
\makeatother
\usepackage{longtable,booktabs,array}
\newcounter{none} % for unnumbered tables
\usepackage{calc} % for calculating minipage widths
% Correct order of tables after \paragraph or \subparagraph
\usepackage{etoolbox}
\makeatletter
\patchcmd\longtable{\par}{\if@noskipsec\mbox{}\fi\par}{}{}
\makeatother
% Allow footnotes in longtable head/foot
\IfFileExists{footnotehyper.sty}{\usepackage{footnotehyper}}{\usepackage{footnote}}
\makesavenoteenv{longtable}
\setlength{\emergencystretch}{3em} % prevent overfull lines
\providecommand{\tightlist}{%
  \setlength{\itemsep}{0pt}\setlength{\parskip}{0pt}}
\usepackage[]{natbib}
\bibliographystyle{plainnat}
\makeatletter
\@ifpackageloaded{subcaption}{}{\usepackage{subcaption}}
\@ifpackageloaded{caption}{}{\usepackage{caption}}
\captionsetup[subfigure]{margin=0.5em}
\AtBeginDocument{%
\renewcommand*\figurename{Figure}
\renewcommand*\tablename{Table}
}
\AtBeginDocument{%
\renewcommand*\listfigurename{List of Figures}
\renewcommand*\listtablename{List of Tables}
}
\newcounter{pandoccrossref@subfigures@footnote@counter}
\newenvironment{pandoccrossrefsubfigures}{%
\setcounter{pandoccrossref@subfigures@footnote@counter}{0}
\begin{figure}\centering%
\gdef\global@pandoccrossref@subfigures@footnotes{}%
\DeclareRobustCommand{\footnote}[1]{\footnotemark%
\stepcounter{pandoccrossref@subfigures@footnote@counter}%
\ifx\global@pandoccrossref@subfigures@footnotes\empty%
\gdef\global@pandoccrossref@subfigures@footnotes{{##1}}%
\else%
\g@addto@macro\global@pandoccrossref@subfigures@footnotes{, {##1}}%
\fi}}%
{\end{figure}%
\addtocounter{footnote}{-\value{pandoccrossref@subfigures@footnote@counter}}
\@for\f:=\global@pandoccrossref@subfigures@footnotes\do{\stepcounter{footnote}\footnotetext{\f}}%
\gdef\global@pandoccrossref@subfigures@footnotes{}}
\@ifpackageloaded{float}{}{\usepackage{float}}
\floatstyle{ruled}
\@ifundefined{c@chapter}{\newfloat{codelisting}{h}{lop}}{\newfloat{codelisting}{h}{lop}[chapter]}
\floatname{codelisting}{Listing}
\newcommand*\listoflistings{\listof{codelisting}{List of Listings}}
\makeatother
\usepackage{bookmark}
\IfFileExists{xurl.sty}{\usepackage{xurl}}{} % add URL line breaks if available
\urlstyle{same}
% fallback for those not using the hyperref driver hyperxmp:
\makeatletter
\@ifundefined{xmpquote}{\newcommand{\xmpquote}[1]{#1}}{}
\makeatother
\hypersetup{
  colorlinks=true,linkcolor=red,urlcolor=red,citecolor=red,anchorcolor=red,filecolor=red,
  pdfcreator={LaTeX via pandoc}}

\author{}
\date{}

\IfFileExists{seqsplit.sty}{\usepackage{seqsplit}}{\newcommand{\seqsplit}[1]{#1}}
\protected\def\breakseq#1{\seqsplit{#1}}
\protected\def\breaktt#1{\begingroup\ttfamily\seqsplit{#1}\endgroup}
\title{The Triplicate: A Data-Driven Large-Format Newspaper Layout Engine}
\author{Daniel Ari Friedman\\\footnotesize{Active Inference Institute}\\\footnotesize{\texttt{daniel@activeinference.institute}}\\\footnotesize{\href{https://orcid.org/0000-0001-6232-9096}{ORCID: 0000-0001-6232-9096}} \\ \footnotesize{DOI: \href{https://doi.org/10.5281/zenodo.20533675}{10.5281/zenodo.20533675}}}
\date{\today}

\usepackage[margin=0.75in]{geometry}
\usepackage{graphicx}

\begin{document}

\begin{titlepage}
\centering
\vspace*{2cm}
{\LARGE\bfseries The Triplicate: A Data-Driven Large-Format Newspaper Layout Engine\par}
\vspace{0.75em}
{\large Designed multi-column page layout as a reproducible-research exemplar\par}
\vfill
\makeatletter
{\@author\par}
\vspace{1em}
{\@date\par}
\makeatother
\vfill
\end{titlepage}
\thispagestyle{empty}
\tableofcontents
\newpage


\section{Abstract}\label{abstract}

We present \breaktt{template\_newspaper}, a pure-Python engine that
renders a complete twelve-page, large-format newspaper to a print-ready
PDF from structured YAML content. The exemplar edition is \emph{The
Triplicate}, a homage to the historic newspaper of Crescent City,
California (founded 1879). The engine demonstrates designed multi-column
page layout --- nameplate and ears, spanning headlines, flowing column
frames with an optional rail, drop caps, pull quotes, ruled modular
boxes, data tables, halftone ``engraving'' figures and running folios
--- while keeping a strict separation between \emph{content} (data) and
\emph{engine} (code): a new title is a data edit, never a code change.
The layout strategy is a hybrid in which fixed \emph{furniture} is drawn
directly on the canvas to establish where the column grid begins, after
which body copy flows through ReportLab frames that split paragraphs
across columns automatically. The project obeys the research-template
monorepo contract and is discovered and executed by the same
orchestration pipeline as its code- and prose-focused siblings.

\newpage

\section{Introduction}\label{introduction}

Most reproducible-research tooling treats a document as a linear stream:
a manuscript is sections of prose, rendered top to bottom. A newspaper
is the opposite --- a deliberately \emph{non-linear}, spatial artifact,
in which a reader's eye is guided by a hierarchy of headline sizes, by
rules and gutters, by the placement of a photograph and the interruption
of a pull quote. Building one programmatically is therefore an unusually
demanding test of a layout system, and an unusually legible
demonstration when it succeeds.

\breaktt{template\_newspaper} exists to be that demonstration within the
research template. It is the third canonical exemplar, joining
\breaktt{template\_code\_project} (a computational pipeline) and
\breaktt{template\_prose\_project} (a manuscript-quality pipeline). All
three share one directory contract and one orchestration pipeline; they
differ only in what they render.

The design follows a single principle: \textbf{content is data, the
engine is the press.} An edition is a small set of YAML files under
\texttt{content/} --- a masthead manifest and one file per page, each
declaring its stories, boxes and figures. The engine in
\texttt{src/newspaper/} reads that data and renders it; it contains no
copy and, outside a single typography module, names no fonts. Producing
a different paper --- a school gazette, a conference daily, a
neighborhood broadsheet --- is an edit to the data, not the code.

The remainder of this manuscript describes the engine's architecture
(Section 2), its typography and figure generation (Section 3), and how
the artifact is produced and verified reproducibly (Section 4).

\newpage

\section{Engine Architecture}\label{engine-architecture}

The engine is nine small modules, each with one responsibility.

\texttt{geometry} is pure arithmetic over points: the page trim, its
margins, and a \texttt{ColumnGrid} that partitions a region into equal
columns with gutters and reports the centre of each gutter (where a
hairline rule is drawn). It imports no rendering library, which makes
the column mathematics trivially testable.

\texttt{content} defines the typed object graph --- \texttt{Edition},
\texttt{Page}, \texttt{Story}, \texttt{Box}, \texttt{Figure},
\texttt{Block} --- and the strict YAML loaders that build it. The
loaders are forgiving about \emph{shape} (a body element may be a bare
string or a tagged mapping) but strict about \emph{structure}: an
unknown page template or box kind raises an error naming the offending
value and the allowed set.

\texttt{config} is the strict render configuration (trim size, default
column count, rail width, gutter), validated the same way.

\texttt{typography} owns every font name: it registers the Didot
display, Georgia text and Helvetica label faces (with a base-14
fall-back when the system lacks them) and builds the paragraph
stylesheet that the rest of the engine draws from through named
\texttt{Fonts} roles.

\texttt{figures} is the only module that \emph{generates} imagery: the
halftone harbour, lighthouse and redwood engravings (Pillow), the
weather, tide and crab-landings charts (Matplotlib), and the colour
display-ad graphics. It writes deterministic PNGs that
\texttt{components} later places.

\texttt{components} builds the flowables that live \emph{inside} a
column: a story's headline, byline and body; a raised drop cap; a pull
quote; a ruled box; a data table; a figure with its cutline.
\texttt{furniture} draws the elements that are \emph{placed} rather than
flowed: the nameplate and its ears, the section banner, the folio, the
hairline column rules, and the spanning lead headline.

\texttt{layout} is the bridge. For each page it draws the furniture
(which returns the y-coordinate where the column grid begins),
constructs the column frames --- including an optional narrow rail and
the room reserved for a spanning headline --- and pours the flowables
into the frames with ReportLab's \breaktt{Frame.addFromList}, which
splits paragraphs across columns automatically. Content that does not
fit is returned and reported as an \emph{over-set} page rather than
silently dropped.

\texttt{engine} ties it together: register fonts, build the stylesheet,
open a canvas at the configured trim, render each page, and write the
PDF plus a machine-readable render report.

This \textbf{drawn-furniture / flowed-body} split is the key
architectural choice. A pure document-template approach cannot place a
headline that spans several columns above body text that flows beneath
it; drawing the furniture first, and beginning the column frames below
it, makes spanning headlines, standing boxes and automatic text flow
coexist on the same page.

\subsection{Methods}\label{methods}

Producing an edition is a deterministic, four-step method, run by the
pipeline and covered end-to-end by the test suite:

\begin{enumerate}
\def\labelenumi{\arabic{enumi}.}
\tightlist
\item
  \textbf{Load} --- \texttt{content} parses \texttt{edition.yaml} and
  every \breaktt{content/pages/*.yaml} into the typed \texttt{Edition}
  graph, rejecting unknown page templates or box kinds by naming the
  offending value and the allowed set.
\item
  \textbf{Measure} --- \texttt{geometry} derives the trim, margins and
  \texttt{ColumnGrid} for each page from the strict \texttt{config},
  independently of any rendering library.
\item
  \textbf{Compose} --- \texttt{components} and \texttt{furniture} build
  the flowed and drawn elements; \texttt{layout} pours them into the
  column frames and reports any content that does not fit as an over-set
  page rather than dropping it silently.
\item
  \textbf{Emit} --- \texttt{engine} registers fonts, renders each page
  to the canvas, and writes the PDF alongside a machine-readable render
  report.
\end{enumerate}

Every step is pure-data-in, artifact-out and seeded deterministically,
so the same edition manifest always yields a byte-identical paper ---
the reproducibility contract verified in \texttt{tests/} and described
in \breaktt{04\_reproducibility.md}.

\newpage

\section{Typography and Figures}\label{typography-and-figures}

A newspaper's voice is its type. The engine registers, when present, the
elegant text faces a real paper would license --- Didot for the
nameplate and display headlines, Georgia for body text, and a grotesque
sans for kickers and labels --- and falls back to ReportLab's
always-available base-14 fonts when a face is missing. Registration is
best-effort and never fatal: a checkout on a machine without the macOS
supplemental fonts still renders a structurally identical paper in Times
and Helvetica. Only the typography module names a font; every other
module refers to four logical roles (display, body, sans, and their bold
and italic variants), so a change of typeface touches one file.

A subtle lesson is recorded in the code: the bold weight inside a
\texttt{.ttc} font collection is at a face-specific subfont index.
Didot's collection orders its faces Regular, Italic, Bold; an early
render produced \emph{italic} headlines because the bold index was
assumed to be one rather than two. The fix --- and the comment that now
guards it --- came from inspecting the rendered raster, not the source.

Figures are generated in the monochrome idiom of newsprint by two means.
Photo- graphic elements --- the harbor, the lighthouse, the redwoods ---
are procedural grayscale scenes drawn with Pillow and passed through a
forty-five-degree halftone screen, so they read as classic dot-screened
newspaper photographs. Data graphics --- a seven-day forecast, a tide
curve, a Dungeness-crab landings bar chart --- are rendered with
Matplotlib in black and gray with serif labels, so they sit naturally
beside the type rather than fighting it. A missing image file degrades
to a clearly labelled placeholder, never a crash.

\newpage

\section{Reproducibility}\label{reproducibility}

The artifact is produced by a deterministic, three-step chain that the
repository pipeline runs as Stage-02 analysis scripts, in lexicographic
order: \texttt{00\_preflight.py} confirms the dependencies import and
the edition loads and validates; \breaktt{10\_generate\_figures.py}
writes the halftone scenes and charts to \texttt{../figures/}; and
\breaktt{20\_render\_newspaper.py} renders the twelve pages to
\breaktt{output/pdf/the-triplicate.pdf} and emits a machine-readable
render report (\breaktt{output/data/render\_report.json}) plus a human
summary. The render is deterministic --- the canvas is opened in
ReportLab's \texttt{invariant} mode and the figures are drawn from fixed
procedural recipes --- so the same content yields the same bytes.

Correctness is established two ways. A \texttt{pytest} suite exercises
the pure logic (geometry, configuration, content loaders, typography,
figure generation) and renders the real edition end-to-end, asserting
that the output is a valid PDF of exactly twelve pages at the correct
trim with no over-set page. Because the remaining correctness of a
\emph{layout} is visual, the development loop rasterizes each page and
inspects it; the project's agent guidance makes that visual check
mandatory after any layout change. The engine is type-checked with mypy
and linted with ruff, and the test suite reports roughly ninety-five
percent line coverage.

Every quantitative claim in this project --- the page count, the trim
dimensions, the figure count --- is recorded in
\breaktt{data/claim\_ledger.yaml} against the artifact that substantiates
it, so a reviewer can trace each number to its source.

\newpage

\section{References}\label{references}

The bibliography is maintained in \texttt{references.bib}. Key sources:

\begin{itemize}
\tightlist
\item
  ReportLab Ltd.~\emph{ReportLab PDF Library --- User Guide.} The
  platypus flowable and frame model underpins this engine's column flow
  \citep{reportlab}.
\item
  Bringhurst, R. \emph{The Elements of Typographic Style.} The measure,
  scale and rule conventions that inform the stylesheet
  \citep{bringhurst}.
\item
  Wong, B. ``Points of view: Color blindness.'' \emph{Nature Methods}
  (2011). The basis for accessible, here monochrome, figure palettes
  \citep{wong2011}.
\item
  García, M. \emph{Contemporary Newspaper Design.} Grid, modular layout
  and the rail/well distinction \citep{garcia}.
\end{itemize}



\clearpage

\bibliography{references}
\end{document}
